\section{Outstanding challenges}
\label{sec:gt-challenges}

The development now supplies several different kinds of information.
The complete error law controls a finite readout; counted contours give
regular pair coordinates; exterior weights convert geometric lower
counts into mass; source density gives averaged and exceptional-pair
estimates; and common-time dynamics improves occupation and interior
interaction bounds. Their hypotheses specify the remaining work.
A bound for an average, a count, or an integrated multiplier is not
interchangeable with the sign of an individual gap.

\subsection{The signed microscopic environment of a selected pair}

The canonical far tail is quantitatively favorable, but at
\(R=\sqrt{\log T}\) its contribution to \(\ell^2M\) tends to zero.
Thus the leading microscopic question lies in the nearby exterior zeros.
The sharp rectangle minimum identifies the best lower weight from its
specified ranges, and real brackets can provide lower counts without
singleton counts or continuously labeled zeros. The supporting finite
application has already acquired the local mass threshold. Repeating
that fixed test is not the missing height-uniform result.

The source mean \(C/a^3\) improves the equally averaged count bound
\(C_0/a\). Negative-incidence estimates also make exceptional pairs
sparse as the normalized collar \(a\) grows. At a fixed \(a\), however,
neither controls a chosen center. The identical-horizontal-count
construction preserves the aggregate moment and density upper bounds
with enlarged constants while changing its selected interaction sign.
A useful further estimate must consume additional conditional or signed
information from the unchanged complete source. It must also acquire,
or explicitly assume, the empty collar to which it applies. An upper
population bound alone supplies neither a favorable lower mass nor a
probability law relative to the possibly sparse eligible-pair population.

\subsection{The forcing in the exact interaction equation}

The equation
\[
 m_\sigma=2m^2-\mathcal F,\qquad
 \mathcal F=12\sum_j\mu_j^2+16\delta\sum_j\mu_j^3
\]
locates a concrete signed quantity beyond the coarse absolute shell
estimate \(|\mathcal F|\le C_F\log T\). Its square and cube sums are
real after conjugate grouping, but need not be termwise nonnegative.
The resulting Riccati comparison improves the normalized negative
allowance to order \(\sqrt{\log T}\), conditional on a sufficiently
long counted tube. The upper bound on its integrated multiplier also
improves, but that upper allowance still grows without bound at fixed
positive scaled duration.

A height-independent or otherwise sharper signed forcing estimate would
have a different quantitative consequence. Merely stating such an
estimate is not progress: its source input, complete remainder and
application domain must be supplied. It may also be necessary to use
joint information about \(D\), \(m\) and the exterior sums rather than
estimating each independently. No favorable sign for the forcing or
incoming gap is proved by the present identity.

\subsection{Time visibility, tube acquisition and incoming data}

Retained conjugate dissipation gives a fixed-collar negative-pair
occupation allowance of order \(T/\log T\) on fixed physical-time
intervals. This does not count distinct events or control a specified
slice. The available physical modulus of \(\ell^2M\) is of order
\((\log T)^4\). Dividing the occupation allowance by its guaranteed
negative-visibility duration loses the aggregate gain. Even granting
negative visibility for a spacing-squared time interval leaves an
allowance of order \(T\log T\), which does not exclude one pair.
These are limitations of the stated upper allowances, not lower bounds
on actual exception counts.

The geometric and dynamical premises must therefore remain separate.
A count-two contour at one time does not automatically persist over a
microscopic time interval. A tube that does persist need not maintain a
negative interaction. A margin supplies visibility only while the tube
and collar are independently controlled. Finally, the regular law
\(D'=2-4MD\) still needs an incoming value or a justified boundary
argument. The identity \(D'=2\) at a double allows passage through zero;
it is not a positive earlier-gap anchor.

\subsection{Joint Fourier information beyond separate phase payments}

Deflated curvature is a two-point quantity; at a stipulated double it
has numerator \(2(F^{(3)})^2-3F''F^{(4)}\) divided by \(18F''{}^2\).
The natural difference-frequency matrix is indefinite when two distinct
positive spectral radii occur, but that matrix statement does not decide
the numerator along its restricted Fourier curve. The complete continuous
source has no zero-frequency atom. Retaining every favorable cosine phase
while paying each discarded phase separately by \(-1\) still gives the
negative limiting allowance
\((1/\pi-1/2)3m_2m_4\). Its positive normalized averages have the same
failure at high centers.

A successful Fourier estimate must therefore retain signed information
that these specified allowances discard, or justify a different exact
grouping. It must pay the actual frequencies and their complete remainder,
then extend from the double slice to the real-segment or vertical-average
quantity appropriate to the pair. The estimates here do not establish
negative canonical curvature, nor do they rule out such genuinely joint
source estimates. They delimit the tested information classes.

\subsection{Uniform families, higher multiplicity and fair comparison}

The critical-scale event results remain conditional on the stipulated
isolated double. The supporting interval certificates establish local
facts on one named tube, without covering new heights or resolving the
canonical zero problem. A global application would require complete
counted coverage, uniform collars and source errors, a justified limit
towards time zero, and composition across the relevant height and time
ranges. Higher-multiplicity events require further symmetric coordinates;
a single squared gap cannot describe their full geometry.

Every comparison must give direct evaluation the same source identities,
exterior counts, coordinate information and incoming data. The local
numerical example already has a stronger direct drift enclosure and an
early-endpoint minimum. The new analytical estimates are shared by an
equally informed direct method; no exclusive transport advantage or total
acquisition-cost saving has been established. The remaining objective is
a quantitative, source-specific estimate that closes one of the identified
conditional or visibility gaps while preserving those common premises.
Global reality and every-pair positivity remain open.
